\chapter{Algorithms --- Sorting Algorithms and Comparative Analysis}
\label{ch:algorithms_sorting}

This chapter covers Bubble Sort, Insertion Sort, Merge Sort, and Quick Sort in detail, followed by other important sorting algorithms frequently encountered in computer science recruitment examinations.

The focus is on algorithm design, step-by-step execution, pseudocode, complete C++ implementations, time and space complexity, stability, in-place sorting, practical applications, and common examination traps.

\begin{center}
\vspace{0.5em}
{\color{primary}\rule{0.45\textwidth}{1.2pt}}\\[4pt]
{\Large\bfseries\color{primary} Section I: Fundamentals of Sorting}\\[2pt]
{\color{primary}\rule{0.45\textwidth}{1.2pt}}
\vspace{0.5em}
\end{center}

\section{Fundamentals of Sorting}

\subsection{What is Sorting?}

\textbf{Sorting} is the process of arranging a collection of data in a specific sequential order according to a designated key or comparison criterion.

The two fundamental orderings are:
\begin{itemize}
    \item \textbf{Ascending order (Non-decreasing):} Elements increase from smallest to largest, e.g., $3, 7, 12, 19$.
    \item \textbf{Descending order (Non-increasing):} Elements decrease from largest to smallest, e.g., $19, 12, 7, 3$.
\end{itemize}

Sorting applies across fundamental data types: integers, floating-point numbers, strings, dates, employee records, database rows, and custom objects. For example, a government recruitment portal sorts applicants by merit marks, roll numbers, or alphabetical name order.

\subsection{Why is Sorting Important?}

Sorting dramatically enhances downstream computational efficiency:
\begin{itemize}
    \item \textbf{Binary Search:} Enables logarithmic $O(\log n)$ item location in arrays with random access.
    \item \textbf{Database Indexing:} B-trees and B+ trees maintain sorted key order to service range queries efficiently.
    \item \textbf{Duplicate Detection:} Identifies duplicate records by scanning contiguous neighbors in $O(n)$ time.
    \item \textbf{Algorithmic Prerequisite:} Greedy algorithms (Kruskal's MST, interval scheduling), divide-and-conquer routines, and convex hull algorithms rely directly on pre-sorted data.
\end{itemize}

\section{Important Sorting Terminology}

Technical recruitment examinations regularly test sorting properties beyond raw asymptotic time complexity:

\subsection{Stable Sorting}

A sorting algorithm is \textbf{stable} if elements with equal sorting keys preserve their original relative order in the sorted output.

\begin{table}[htbp]
\centering
\small
\begin{tabular}{c l c}
\toprule
\textbf{Original Index} & \textbf{Candidate Name} & \textbf{Exam Marks} \\
\midrule
1 & Rahim  & 85 \\
2 & Karim  & 92 \\
3 & Salma  & 85 \\
4 & Nabila & 78 \\
\bottomrule
\end{tabular}
\caption{Candidate records prior to sorting by marks.}
\end{table}

Sorting by marks in ascending order:
\begin{itemize}
    \item \textbf{Stable Result:} $\text{Nabila (78)} \to \textbf{Rahim (85)} \to \textbf{Salma (85)} \to \text{Karim (92)}$. Rahim appears before Salma because Rahim preceded Salma in the input.
    \item \textbf{Unstable Result:} $\text{Nabila (78)} \to \textbf{Salma (85)} \to \textbf{Rahim (85)} \to \text{Karim (92)}$. Salma is placed ahead of Rahim, reversing their input chronology.
\end{itemize}

\begin{quote}
\textbf{Engineering Significance:} Stability is vital in multi-key sorting (e.g., sorting first by Student Name, then stably sorting by Department ID, preserving alphabetical order within each department).
\end{quote}

\subsection{In-Place Sorting}

An algorithm is \textbf{in-place} if it transforms the input using $O(1)$ (constant) auxiliary memory beyond the original input container, allowing at most $O(\log n)$ stack space for recursion:
\begin{itemize}
    \item \textbf{In-place algorithms:} Bubble Sort, Insertion Sort, Selection Sort, Heap Sort, and array-based Quick Sort.
    \item \textbf{Not in-place algorithms:} Conventional array Merge Sort ($O(n)$ auxiliary array), Counting Sort, and Radix Sort.
\end{itemize}

\subsection{Adaptive Sorting}

An algorithm is \textbf{adaptive} if its execution time improves when the input sequence is already sorted or nearly sorted:
\begin{itemize}
    \item \textbf{Insertion Sort} is inherently adaptive, running in $O(n)$ linear time on already sorted inputs.
    \item \textbf{Optimized Bubble Sort} (with a swap-detection flag) is adaptive ($O(n)$ on sorted input).
    \item \textbf{Merge Sort} is non-adaptive in its standard top-down implementation, performing $\Theta(n \log n)$ operations regardless of initial order.
\end{itemize}

\subsection{Comparison-Based vs Non-Comparison-Based Sorting}

\begin{itemize}
    \item \textbf{Comparison-Based Sorts:} Determine relative order exclusively by pairwise comparisons ($A[i] \leq A[j]$). Examples: Bubble, Insertion, Selection, Merge, Quick, and Heap Sort. \textbf{Theoretical Lower Bound:} $\Omega(n \log n)$ worst-case time.
    \item \textbf{Non-Comparison-Based Sorts:} Exploit intrinsic mathematical properties or digit representations of keys (e.g., bounded integer ranges). Examples: Counting Sort, Radix Sort, Bucket Sort. Can achieve linear $O(n)$ time.
\end{itemize}

\subsection{Internal vs External Sorting}

\begin{itemize}
    \item \textbf{Internal Sorting:} All data resides entirely inside primary random-access memory (RAM).
    \item \textbf{External Sorting:} The dataset exceeds RAM capacity and resides on secondary external storage (SSD/HDD). Algorithms (e.g., External Merge Sort) optimize sequential block I/O.
\end{itemize}

\newpage
\begin{center}
\vspace{0.5em}
{\color{primary}\rule{0.45\textwidth}{1.2pt}}\\[4pt]
{\Large\bfseries\color{primary} Section II: Bubble Sort}\\[2pt]
{\color{primary}\rule{0.45\textwidth}{1.2pt}}
\vspace{0.5em}
\end{center}

\section{Bubble Sort}

\textbf{Bubble Sort} is a simple comparison-based sorting algorithm that repeatedly steps through the array, compares adjacent elements, and swaps them if they are in the wrong order. With each complete pass, the largest unsorted element "bubbles up" to its correct position at the end of the array.

\subsection{Algorithmic Steps}
\begin{enumerate}
    \item Traverse the array from index $0$ to $n-2$.
    \item Compare adjacent pair $(A[j], A[j+1])$. If $A[j] > A[j+1]$, swap them.
    \item At the end of pass $i$, the $(i+1)^{\text{th}}$ largest element is locked into position $n - 1 - i$.
    \item If a complete pass finishes without a single swap, terminate early (array is sorted).
\end{enumerate}

\subsection{Step-by-Step Example}
Sort $A = [5, 1, 4, 2, 8]$ in ascending order:
\begin{itemize}
    \item \textbf{Pass 1:}
    \begin{itemize}
        \item Compare $5, 1 \implies \text{swap} \implies [1, 5, 4, 2, 8]$
        \item Compare $5, 4 \implies \text{swap} \implies [1, 4, 5, 2, 8]$
        \item Compare $5, 2 \implies \text{swap} \implies [1, 4, 2, 5, 8]$
        \item Compare $5, 8 \implies \text{no swap} \implies [1, 4, 2, 5, 8]$
        \item Largest element $8$ is in its permanent position.
    \end{itemize}
    \item \textbf{Pass 2:}
    \begin{itemize}
        \item Compare $1, 4 \implies \text{no swap} \implies [1, 4, 2, 5, 8]$
        \item Compare $4, 2 \implies \text{swap} \implies [1, 2, 4, 5, 8]$
        \item Compare $4, 5 \implies \text{no swap} \implies [1, 2, 4, 5, 8]$
        \item Second largest element $5$ is locked into place.
    \end{itemize}
    \item \textbf{Pass 3:}
    \begin{itemize}
        \item Compare $1, 2 \implies \text{no swap}$; Compare $2, 4 \implies \text{no swap}$.
        \item No swaps occurred $\implies$ \textbf{Early termination triggered!}
    \end{itemize}
\end{itemize}
Final sorted array: $\mathbf{[1, 2, 4, 5, 8]}$.

\subsection{Pseudocode}

\begin{lstlisting}[style=pseudocode, caption={Optimized Bubble Sort Pseudocode.}]
BUBBLE_SORT(A, n):
    for i = 0 to n - 2:
        swapped = false
        for j = 0 to n - 2 - i:
            if A[j] > A[j + 1]:
                swap(A[j], A[j + 1])
                swapped = true
        if swapped == false:
            break
\end{lstlisting}

\subsection{C++ Implementation}

\begin{lstlisting}[style=cppstyle, caption={C++ Optimized Bubble Sort Implementation.}]
#include <iostream>
#include <vector>
#include <utility>
using namespace std;

void bubbleSort(vector<int>& a) {
    int n = static_cast<int>(a.size());
    for (int i = 0; i < n - 1; i++) {
        bool swapped = false;
        for (int j = 0; j < n - 1 - i; j++) {
            if (a[j] > a[j + 1]) {
                swap(a[j], a[j + 1]);
                swapped = true;
            }
        }
        if (!swapped) break; // Early termination
    }
}

int main() {
    vector<int> a = {5, 1, 4, 2, 8};
    bubbleSort(a);
    for (int x : a) cout << x << " ";
    cout << '\n';
    return 0;
}
\end{lstlisting}

\subsection{Complexity Analysis}
\begin{itemize}
    \item \textbf{Worst-Case Time:} Occurs when the array is in reverse sorted order. Total comparisons:
    \[
    (n-1) + (n-2) + \dots + 1 = \frac{n(n-1)}{2} \implies \boxed{O(n^2)}
    \]
    \item \textbf{Average-Case Time:} Quadratic pairwise inversions $\implies \boxed{O(n^2)}$.
    \item \textbf{Best-Case Time:} Occurs when array is already sorted. The \texttt{swapped} flag triggers exit after $n-1$ comparisons $\implies \boxed{O(n)}$.
    \item \textbf{Auxiliary Space:} Modifies the array in place $\implies \boxed{O(1)}$.
\end{itemize}

\subsection{Properties Summary}

\begin{table}[htbp]
\centering
\small
\begin{tabularx}{\textwidth}{l X}
\toprule
\textbf{Property} & \textbf{Characteristics} \\
\midrule
Best-Case Time    & $O(n)$ (with \texttt{swapped} flag optimization) \\
Average-Case Time & $O(n^2)$ \\
Worst-Case Time   & $O(n^2)$ \\
Auxiliary Space   & $O(1)$ (In-place) \\
Stability         & \textbf{Stable} (equal elements are never swapped when using $A[j] > A[j+1]$) \\
Adaptivity        & \textbf{Adaptive} (benefits from existing order) \\
Practical Usage   & Primarily educational; useful for tiny or nearly sorted lists \\
\bottomrule
\end{tabularx}
\caption{Properties of Bubble Sort.}
\end{table}

\newpage
\begin{center}
\vspace{0.5em}
{\color{primary}\rule{0.45\textwidth}{1.2pt}}\\[4pt]
{\Large\bfseries\color{primary} Section III: Insertion Sort}\\[2pt]
{\color{primary}\rule{0.45\textwidth}{1.2pt}}
\vspace{0.5em}
\end{center}

\section{Insertion Sort}

\textbf{Insertion Sort} incrementally constructs a sorted subarray at the front of the array. For each new element (the \texttt{key}), it scans backward through the sorted prefix, shifts larger elements one position to the right, and inserts the key into its correct relative position.

\subsection{Algorithmic Steps}
\begin{enumerate}
    \item The first element $A[0]$ is trivially sorted.
    \item For $i = 1$ to $n-1$: designate $\text{key} = A[i]$ and set $j = i - 1$.
    \item While $j \geq 0$ and $A[j] > \text{key}$: shift $A[j+1] = A[j]$, and decrement $j$.
    \item Place $\text{key}$ at the vacated slot: $A[j+1] = \text{key}$.
\end{enumerate}

\subsection{Step-by-Step Example}
Sort $A = [5, 1, 4, 2, 8]$:
\begin{itemize}
    \item Initial sorted prefix: $[5]$.
    \item \textbf{Pass 1 ($i = 1$, key = 1):} $5 > 1 \implies$ shift $5 \implies$ insert $1 \implies \mathbf{[1, 5}, 4, 2, 8]$.
    \item \textbf{Pass 2 ($i = 2$, key = 4):} $5 > 4 \implies$ shift $5$; $1 < 4 \implies$ stop $\implies$ insert $4 \implies \mathbf{[1, 4, 5}, 2, 8]$.
    \item \textbf{Pass 3 ($i = 3$, key = 2):} $5 > 2, 4 > 2 \implies$ shift $5, 4$; $1 < 2 \implies$ stop $\implies$ insert $2 \implies \mathbf{[1, 2, 4, 5}, 8]$.
    \item \textbf{Pass 4 ($i = 4$, key = 8):} $8 > 5 \implies$ no shifts required $\implies \mathbf{[1, 2, 4, 5, 8]}$.
\end{itemize}

\subsection{Pseudocode}

\begin{lstlisting}[style=pseudocode, caption={Insertion Sort Pseudocode.}]
INSERTION_SORT(A, n):
    for i = 1 to n - 1:
        key = A[i]
        j = i - 1
        while j >= 0 and A[j] > key:
            A[j + 1] = A[j]
            j = j - 1
        A[j + 1] = key
\end{lstlisting}

\begin{quote}
\textbf{Stability Condition:} Using strict inequality \texttt{A[j] > key} ensures stability. If non-strict \texttt{A[j] >= key} were used, equal keys would shift past each other, destroying stability!
\end{quote}

\subsection{C++ Implementation}

\begin{lstlisting}[style=cppstyle, caption={C++ Insertion Sort Implementation.}]
#include <iostream>
#include <vector>
using namespace std;

void insertionSort(vector<int>& a) {
    int n = static_cast<int>(a.size());
    for (int i = 1; i < n; i++) {
        int key = a[i];
        int j = i - 1;
        while (j >= 0 && a[j] > key) {
            a[j + 1] = a[j];
            j--;
        }
        a[j + 1] = key;
    }
}

int main() {
    vector<int> a = {5, 1, 4, 2, 8};
    insertionSort(a);
    for (int x : a) cout << x << " ";
    cout << '\n';
    return 0;
}
\end{lstlisting}

\subsection{Complexity and Properties}

\begin{table}[htbp]
\centering
\small
\begin{tabularx}{\textwidth}{l X}
\toprule
\textbf{Property} & \textbf{Characteristics} \\
\midrule
Best-Case Time    & $\boxed{O(n)}$ (Array already sorted; inner loop test fails immediately) \\
Average-Case Time & $\boxed{O(n^2)}$ (Average $n^2/4$ shifts) \\
Worst-Case Time   & $\boxed{O(n^2)}$ (Reverse sorted array; $\frac{n(n-1)}{2}$ shifts) \\
Auxiliary Space   & $\boxed{O(1)}$ (In-place) \\
Stability         & \textbf{Stable} \\
Online Property   & \textbf{Online} (can sort a live stream of data as elements arrive) \\
Practical Usage   & Base sort for small subarrays ($n \leq 16$--$32$) in hybrid algorithms (TimSort, Introsort) \\
\bottomrule
\end{tabularx}
\caption{Properties of Insertion Sort.}
\end{table}

\newpage
\begin{center}
\vspace{0.5em}
{\color{primary}\rule{0.45\textwidth}{1.2pt}}\\[4pt]
{\Large\bfseries\color{primary} Section IV: Merge Sort}\\[2pt]
{\color{primary}\rule{0.45\textwidth}{1.2pt}}
\vspace{0.5em}
\end{center}

\section{Merge Sort}

\textbf{Merge Sort} is an optimal divide-and-conquer sorting algorithm. It recursively bisects the array into halves until subarrays contain single elements, and then merges the sorted halves into a complete sorted array in linear time.

\subsection{Algorithmic Structure}
\begin{enumerate}
    \item \textbf{Divide:} Bisect the array into two halves at midpoint $\text{mid} = \text{left} + (\text{right} - \text{left})/2$.
    \item \textbf{Conquer:} Recursively invoke \texttt{mergeSort} on left half $[\text{left}, \text{mid}]$ and right half $[\text{mid}+1, \text{right}]$.
    \item \textbf{Combine:} Merge the two sorted subarrays using a linear two-pointer merge procedure.
\end{enumerate}

\subsection{Step-by-Step Example}
Sort $A = [5, 1, 4, 2, 8]$:
\begin{enumerate}
    \item \textbf{Divide Tree:}
    \[
    [5, 1, 4, 2, 8] \implies [5, 1] \text{ and } [4, 2, 8]
    \]
    \[
    [5, 1] \implies [5], [1]; \qquad [4, 2, 8] \implies [4], [2, 8] \implies [2], [8]
    \]
    \item \textbf{Merge Stage 1:}
    \begin{itemize}
        \item Merge $[5]$ and $[1] \implies [1, 5]$.
        \item Merge $[2]$ and $[8] \implies [2, 8]$.
    \end{itemize}
    \item \textbf{Merge Stage 2:}
    \begin{itemize}
        \item Merge $[4]$ with $[2, 8] \implies [2, 4, 8]$.
    \end{itemize}
    \item \textbf{Merge Stage 3 (Final Merge):}
    \begin{itemize}
        \item Merge $[1, 5]$ and $[2, 4, 8]$:
        \item Compare $1 < 2 \implies \text{take } 1$
        \item Compare $5 > 2 \implies \text{take } 2$
        \item Compare $5 > 4 \implies \text{take } 4$
        \item Compare $5 < 8 \implies \text{take } 5$
        \item Append remainder $8 \implies \mathbf{[1, 2, 4, 5, 8]}$.
    \end{itemize}
\end{enumerate}

\subsection{C++ Implementation}

\begin{lstlisting}[style=cppstyle, caption={C++ Merge Sort Implementation.}]
#include <iostream>
#include <vector>
using namespace std;

void merge(vector<int>& a, int left, int mid, int right) {
    vector<int> temp;
    temp.reserve(right - left + 1);

    int i = left;
    int j = mid + 1;

    while (i <= mid && j <= right) {
        if (a[i] <= a[j]) { // '<=' ensures stability
            temp.push_back(a[i++]);
        } else {
            temp.push_back(a[j++]);
        }
    }

    while (i <= mid)   temp.push_back(a[i++]);
    while (j <= right) temp.push_back(a[j++]);

    for (int k = 0; k < static_cast<int>(temp.size()); k++) {
        a[left + k] = temp[k];
    }
}

void mergeSort(vector<int>& a, int left, int right) {
    if (left >= right) return;
    int mid = left + (right - left) / 2;
    mergeSort(a, left, mid);
    mergeSort(a, mid + 1, right);
    merge(a, left, mid, right);
}

int main() {
    vector<int> a = {5, 1, 4, 2, 8};
    mergeSort(a, 0, static_cast<int>(a.size()) - 1);
    for (int x : a) cout << x << " ";
    cout << '\n';
    return 0;
}
\end{lstlisting}

\subsection{Complexity Analysis}
\begin{itemize}
    \item \textbf{Recurrence Relation:}
    \[
    \boxed{T(n) = 2T(n/2) + O(n)}
    \]
    By the Master Theorem (Case 2: $a=2, b=2, f(n)=n \implies n^{\log_2 2} = n$), the solution is:
    \[
    \boxed{T(n) = O(n \log n)}
    \]
    This bound holds strictly across \textbf{Best, Average, and Worst Cases}.
    \item \textbf{Space Complexity:} The merge procedure requires an auxiliary buffer of size $n$, yielding $\boxed{O(n)}$ auxiliary space for array sorting.
\end{itemize}

\subsection{Properties and Architectural Strengths}

\begin{table}[htbp]
\centering
\small
\begin{tabularx}{\textwidth}{l X}
\toprule
\textbf{Property} & \textbf{Characteristics} \\
\midrule
Best / Average / Worst Time & $\boxed{O(n \log n)}$ guaranteed \\
Auxiliary Space             & $O(n)$ for arrays; $O(\log n)$ call stack \\
Stability                   & \textbf{Stable} (preserves order when $A[i] \leq A[j]$) \\
In-Place?                   & \textbf{No} (requires $O(n)$ temporary buffer for arrays) \\
Linked Lists Suitability    & \textbf{Ideal}: Linked lists can be merged in $O(1)$ extra space by pointer redirection without array shifting! \\
External Sorting            & \textbf{Standard choice}: Excellent sequential streaming access patterns on secondary disks. \\
\bottomrule
\end{tabularx}
\caption{Properties of Merge Sort.}
\end{table}

\newpage
\begin{center}
\vspace{0.5em}
{\color{primary}\rule{0.45\textwidth}{1.2pt}}\\[4pt]
{\Large\bfseries\color{primary} Section V: Quick Sort}\\[2pt]
{\color{primary}\rule{0.45\textwidth}{1.2pt}}
\vspace{0.5em}
\end{center}

\section{Quick Sort}

\textbf{Quick Sort} is a divide-and-conquer algorithm that selects a \textbf{pivot element} and partitions the array such that elements smaller than the pivot are placed to its left, and elements greater are placed to its right. The pivot is locked into its final sorted position, and the partitions are sorted recursively.

\subsection{Pivot Selection Strategies}
\begin{itemize}
    \item \textbf{Last element / First element:} Simplest, but degrades to $O(n^2)$ on already sorted arrays.
    \item \textbf{Randomized Pivot:} Selects a random index; guarantees $O(n \log n)$ expected time regardless of input ordering.
    \item \textbf{Median-of-Three:} Selects the median of $\{\text{first}, \text{middle}, \text{last}\}$ elements.
\end{itemize}

\subsection{Lomuto Partition Scheme}
Using the last element as pivot:
\begin{lstlisting}[style=pseudocode, caption={Lomuto Partition Algorithm.}]
PARTITION(A, low, high):
    pivot = A[high]
    i = low - 1
    for j = low to high - 1:
        if A[j] <= pivot:
            i = i + 1
            swap(A[i], A[j])
    swap(A[i + 1], A[high])
    return i + 1
\end{lstlisting}

\subsection{C++ Implementation}

\begin{lstlisting}[style=cppstyle, caption={C++ Quick Sort Implementation.}]
#include <iostream>
#include <vector>
#include <utility>
using namespace std;

int partitionArray(vector<int>& a, int low, int high) {
    int pivot = a[high];
    int i = low - 1;
    for (int j = low; j < high; j++) {
        if (a[j] <= pivot) {
            i++;
            swap(a[i], a[j]);
        }
    }
    swap(a[i + 1], a[high]);
    return i + 1;
}

void quickSort(vector<int>& a, int low, int high) {
    if (low >= high) return;
    int pivotIndex = partitionArray(a, low, high);
    quickSort(a, low, pivotIndex - 1);
    quickSort(a, pivotIndex + 1, high);
}

int main() {
    vector<int> a = {5, 1, 4, 2, 8};
    quickSort(a, 0, static_cast<int>(a.size()) - 1);
    for (int x : a) cout << x << " ";
    cout << '\n';
    return 0;
}
\end{lstlisting}

\subsection{Complexity Analysis}
\begin{itemize}
    \item \textbf{Best Case ($O(n \log n)$):} Pivot divides array into balanced halves ($n/2, n/2$):
    \[
    T(n) = 2T(n/2) + O(n) \implies O(n \log n)
    \]
    \item \textbf{Average Case ($O(n \log n)$):} Even unbalanced splits (e.g., $90\% / 10\%$) yield logarithmic tree depth $\implies O(n \log n)$.
    \item \textbf{Worst Case ($O(n^2)$):} Occurs when pivot is always the smallest or largest element (e.g., already sorted array with last-element pivot):
    \[
    T(n) = T(n-1) + O(n) \implies O(n^2)
    \]
    \item \textbf{Space Complexity:} In-place array modification, but recursion stack uses:
    \[
    \boxed{O(\log n) \text{ average}, \quad O(n) \text{ worst-case}}
    \]
\end{itemize}

\subsection{Properties Summary}

\begin{table}[htbp]
\centering
\small
\begin{tabularx}{\textwidth}{l X}
\toprule
\textbf{Property} & \textbf{Characteristics} \\
\midrule
Best-Case Time    & $O(n \log n)$ \\
Average-Case Time & $O(n \log n)$ \\
Worst-Case Time   & $\boxed{O(n^2)}$ \\
Auxiliary Space   & $O(\log n)$ average recursion stack; $O(n)$ worst-case \\
Stability         & \textbf{Unstable} (long-distance swaps disrupt relative order of equal keys) \\
In-Place?         & \textbf{Yes} (partitions array in place) \\
Performance Note  & Best cache locality and lowest constant factors among $O(n \log n)$ comparison sorts \\
\bottomrule
\end{tabularx}
\caption{Properties of Quick Sort.}
\end{table}

\newpage
\begin{center}
\vspace{0.5em}
{\color{primary}\rule{0.45\textwidth}{1.2pt}}\\[4pt]
{\Large\bfseries\color{primary} Section VI: Comparative Analysis of Core Algorithms}\\[2pt]
{\color{primary}\rule{0.45\textwidth}{1.2pt}}
\vspace{0.5em}
\end{center}

\section{Comparative Analysis of Core Algorithms}

\subsection{Master Evaluation of the Four Fundamental Sorts}

\begin{table}[htbp]
\centering
\footnotesize
\begin{tabularx}{\textwidth}{p{2.6cm} X X X X}
\toprule
\textbf{Feature} & \textbf{Bubble Sort} & \textbf{Insertion Sort} & \textbf{Merge Sort} & \textbf{Quick Sort} \\
\midrule
Paradigm & Adjacent comparisons & Incremental sorted prefix insertion & Divide-and-Conquer & Partition around Pivot \\
Best Time & $O(n)$ (optimized) & $O(n)$ & $O(n \log n)$ & $O(n \log n)$ \\
Average Time & $O(n^2)$ & $O(n^2)$ & $O(n \log n)$ & $O(n \log n)$ \\
Worst Time & $O(n^2)$ & $O(n^2)$ & $O(n \log n)$ & $O(n^2)$ \\
Space & $O(1)$ & $O(1)$ & $O(n)$ (Buffer) & $O(\log n)$ (Stack) \\
Stability & \textbf{Stable} & \textbf{Stable} & \textbf{Stable} & \textbf{Unstable} \\
In-Place & \textbf{Yes} & \textbf{Yes} & \textbf{No} & \textbf{Yes} \\
Adaptive & \textbf{Yes} & \textbf{Yes} & \textbf{No} & \textbf{No} \\
Primary Strength & Simple to teach & Fastest on small / nearly sorted data & Predictable $O(n \log n)$ and stable & High cache speed, in-place \\
Primary Weakness & Sluggish $O(n^2)$ & Sluggish $O(n^2)$ on large data & Requires $O(n)$ extra memory & Unstable; $O(n^2)$ worst case \\
\bottomrule
\end{tabularx}
\caption{Comprehensive comparison of the four fundamental sorting algorithms.}
\end{table}

\subsection{Strategic Decision Criteria}
\begin{itemize}
    \item \textbf{Merge Sort vs Quick Sort:} Merge Sort guarantees $O(n \log n)$ and stability at the cost of $O(n)$ memory. Quick Sort avoids extra memory and provides superior cache performance, but can degrade to $O(n^2)$ and is unstable.
    \item \textbf{Insertion Sort vs Bubble Sort:} Both are $O(n^2)$ average and $O(n)$ best, but Insertion Sort performs far fewer writes and comparisons, making it universally preferred for small arrays.
\end{itemize}

\newpage
\begin{center}
\vspace{0.5em}
{\color{primary}\rule{0.45\textwidth}{1.2pt}}\\[4pt]
{\Large\bfseries\color{primary} Section VII: Other Important Sorting Algorithms}\\[2pt]
{\color{primary}\rule{0.45\textwidth}{1.2pt}}
\vspace{0.5em}
\end{center}

\section{Other Important Sorting Algorithms}

\subsection{Selection Sort}
Repeatedly identifies the minimum element from the unsorted section and swaps it with the first unsorted element.
\begin{itemize}
    \item \textbf{Time Complexity:} $\boxed{\Theta(n^2)}$ in Best, Average, and Worst cases (always performs $\frac{n(n-1)}{2}$ comparisons).
    \item \textbf{Auxiliary Space:} $O(1)$ (In-place).
    \item \textbf{Stability:} \textbf{Unstable} (e.g., sorting $[4a, 4b, 1]$ swaps $4a$ with $1$, placing $4a$ after $4b$).
    \item \textbf{Key Feature:} Performs at most $\boxed{n - 1 \text{ swaps}}$, making it optimal when \textbf{memory write operations are exceptionally expensive}.
\end{itemize}

\subsection{Heap Sort}
Builds a max-heap in $O(n)$ time, then repeatedly extracts the root maximum and sifts down:
\begin{itemize}
    \item \textbf{Time Complexity:} $\boxed{O(n \log n)}$ in Best, Average, and Worst cases.
    \item \textbf{Auxiliary Space:} $\boxed{O(1)}$ (In-place array manipulation).
    \item \textbf{Stability:} \textbf{Unstable}.
    \item \textbf{Comparison with Quick Sort:} Guarantees $O(n \log n)$ without extra memory, but exhibits poorer hardware cache locality due to jumping indices ($2i + 1, 2i + 2$).
\end{itemize}

\subsection{Shell Sort}
Diminishing increment sort extending Insertion Sort by comparing elements separated by a \textbf{gap} $h$, shrinking $h$ until $h = 1$:
\begin{itemize}
    \item \textbf{Complexity:} Strongly depends on the gap sequence (e.g., Shell's original $N/2^k$ is $O(n^2)$; Knuth's $(3^k - 1)/2$ is $O(n^{3/2})$; Sedgewick's sequence achieves $O(n^{4/3})$).
    \item \textbf{Space:} $O(1)$ in-place; \textbf{Unstable}.
\end{itemize}

\subsection{Counting Sort}
Non-comparison integer sorting that counts occurrences of each value in range $[0, k]$:
\begin{itemize}
    \item \textbf{Time Complexity:} $\boxed{O(n + k)}$. Linear when $k = O(n)$.
    \item \textbf{Auxiliary Space:} $\boxed{O(n + k)}$ using frequency count and output arrays.
    \item \textbf{Stability:} \textbf{Stable} (when iterating input backward during placement).
    \item \textbf{Constraint:} Inefficient when $k \gg n$ (e.g., sorting 10 keys in range $1$ to $1,000,000$).
\end{itemize}

\subsection{Radix Sort}
Processes keys digit-by-digit from Least Significant Digit (LSD) to Most Significant Digit (MSD) using a stable sub-sort (typically Counting Sort):
\begin{itemize}
    \item \textbf{Time Complexity:} $\boxed{O(d(n + b))}$, where $d$ is digit count and $b$ is base ($b = 10$).
    \item \textbf{Auxiliary Space:} $O(n + b)$; \textbf{Stable}.
    \item \textbf{Linear Time:} For fixed-width 32-bit or 64-bit integers, $d$ is constant, yielding linear $O(n)$ sorting!
\end{itemize}

\subsection{Bucket Sort}
Distributes $n$ uniformly distributed elements over $[0, 1)$ into $k$ buckets using index $\lfloor n \cdot x \rfloor$, sorts individual buckets, and concatenates them:
\begin{itemize}
    \item \textbf{Time Complexity:} Expected $\boxed{O(n + k)}$ on uniform inputs; Worst-case $\boxed{O(n^2)}$ under heavy clustering.
    \item \textbf{Space:} $O(n + k)$.
\end{itemize}

\subsection{TimSort}
State-of-the-art hybrid algorithm combining Insertion Sort and Merge Sort:
\begin{itemize}
    \item Identifies natural contiguous ascending/descending subsequences (\textbf{runs}).
    \item Extends short runs using Insertion Sort, then merges them with galloping mode.
    \item \textbf{Time Complexity:} Best $O(n)$ (strictly sorted data); Average/Worst $\boxed{O(n \log n)}$.
    \item \textbf{Space:} $O(n)$; \textbf{Stable}.
    \item \textbf{Standard Adoption:} Default sorting engine in Python (\texttt{list.sort}), Java (\texttt{Arrays.sort} for objects), and Android.
\end{itemize}

\subsection{Introsort (Introspective Sort)}
Hybrid algorithm providing Quick Sort's speed with guaranteed $O(n \log n)$ worst-case safety:
\begin{itemize}
    \item Begins with Quick Sort.
    \item Monitors recursion depth; if depth exceeds $2 \lfloor \log_2 n \rfloor$, switches automatically to \textbf{Heap Sort}.
    \item Switches to \textbf{Insertion Sort} for tiny partitions ($n < 16$).
    \item \textbf{Time Complexity:} Guaranteed $\boxed{O(n \log n)}$ worst-case; Space $O(\log n)$; \textbf{Unstable}.
    \item \textbf{Standard Adoption:} Standard implementation of C++ \texttt{std::sort()}.
\end{itemize}

\subsection{Tree Sort}
Inserts all $n$ items into a Binary Search Tree (BST), then traverses it \textbf{inorder}:
\begin{itemize}
    \item \textbf{Time Complexity:} $O(n \log n)$ average; degrades to $\boxed{O(n^2)}$ on skewed trees. With self-balancing trees (AVL/Red-Black), guarantees $O(n \log n)$.
    \item \textbf{Space:} $O(n)$ for tree nodes.
\end{itemize}

\subsection{External Merge Sort}
Designed for massive datasets exceeding physical RAM:
\begin{enumerate}
    \item \textbf{Run Generation:} Stream chunks of data into RAM, sort internally, and write sorted runs to disk.
    \item \textbf{Multi-Way Merge:} Concurrently merge $k$ sorted runs into a unified output stream using min-heaps.
\end{enumerate}

\newpage
\begin{center}
\vspace{0.5em}
{\color{primary}\rule{0.45\textwidth}{1.2pt}}\\[4pt]
{\Large\bfseries\color{primary} Section VIII: Master Comparative Reference Table}\\[2pt]
{\color{primary}\rule{0.45\textwidth}{1.2pt}}
\vspace{0.5em}
\end{center}

\section{Master Comparative Reference Table}

\begin{table}[htbp]
\centering
\footnotesize
\begin{tabularx}{\textwidth}{p{3.2cm} p{1.8cm} p{1.8cm} p{1.8cm} p{2.6cm} c c}
\toprule
\textbf{Algorithm} & \textbf{Best Time} & \textbf{Average Time} & \textbf{Worst Time} & \textbf{Auxiliary Space} & \textbf{Stable?} & \textbf{In-Place?} \\
\midrule
Bubble Sort (Opt) & $O(n)$ & $O(n^2)$ & $O(n^2)$ & $O(1)$ & Yes & Yes \\
Insertion Sort    & $O(n)$ & $O(n^2)$ & $O(n^2)$ & $O(1)$ & Yes & Yes \\
Selection Sort    & $O(n^2)$ & $O(n^2)$ & $O(n^2)$ & $O(1)$ & No & Yes \\
Merge Sort        & $O(n \log n)$ & $O(n \log n)$ & $O(n \log n)$ & $O(n)$ & Yes & No \\
Quick Sort        & $O(n \log n)$ & $O(n \log n)$ & $O(n^2)$ & $O(\log n)$ stack & No & Yes \\
Heap Sort         & $O(n \log n)$ & $O(n \log n)$ & $O(n \log n)$ & $O(1)$ & No & Yes \\
Shell Sort        & $O(n \log n)$ & $O(n^{4/3})$ & $O(n^2)$ & $O(1)$ & No & Yes \\
Counting Sort     & $O(n + k)$ & $O(n + k)$ & $O(n + k)$ & $O(n + k)$ & Yes & No \\
Radix Sort        & $O(d(n + b))$ & $O(d(n + b))$ & $O(d(n + b))$ & $O(n + b)$ & Yes & No \\
Bucket Sort       & $O(n + k)$ & $O(n + k)$ & $O(n^2)$ & $O(n + k)$ & Yes & No \\
TimSort           & $O(n)$ & $O(n \log n)$ & $O(n \log n)$ & $O(n)$ & Yes & No \\
Introsort         & $O(n \log n)$ & $O(n \log n)$ & $O(n \log n)$ & $O(\log n)$ stack & No & Yes \\
Tree Sort (BST)   & $O(n \log n)$ & $O(n \log n)$ & $O(n^2)$ & $O(n)$ & No & No \\
Pigeonhole Sort   & $O(n + k)$ & $O(n + k)$ & $O(n + k)$ & $O(n + k)$ & Yes & No \\
Comb Sort         & $O(n \log n)$ & $O(n^2 / 2^p)$ & $O(n^2)$ & $O(1)$ & No & Yes \\
\bottomrule
\end{tabularx}
\caption{Master comparative table of 15 sorting algorithms.}
\end{table}

\subsection{Fundamental Comparison-Sorting Theorems}

\begin{tcolorbox}[title=\textbf{Theorem: Lower Bound for Comparison-Based Sorting}]
Any comparison-based sorting algorithm requires at least $\mathbf{\Omega(n \log n)}$ comparisons in the worst case.\\
\textbf{Proof Sketch:} A decision tree sorting $n$ distinct keys must have at least $n!$ leaves (one for each permutation). A binary decision tree of height $h$ has at most $2^h$ leaves:
\[
2^h \geq n! \implies h \geq \log_2(n!)
\]
By Stirling's approximation ($\ln n! \approx n \ln n - n$):
\[
h \geq \sum_{i=1}^n \log_2 i \geq \sum_{i=n/2}^n \log_2(n/2) = \frac{n}{2} \log_2\left(\frac{n}{2}\right) = \Omega(n \log n).
\]
\end{tcolorbox}

\newpage
\begin{center}
\vspace{0.5em}
{\color{primary}\rule{0.45\textwidth}{1.2pt}}\\[4pt]
{\Large\bfseries\color{primary} Section IX: Solved Examination Problems}\\[2pt]
{\color{primary}\rule{0.45\textwidth}{1.2pt}}
\vspace{0.5em}
\end{center}

\section{Solved Examination Problems}

\subsection{Problem 1: Trace Bubble Sort Swaps}
\textbf{Problem Statement:} Sort $A = [7, 3, 5, 2]$ using Bubble Sort. Determine intermediate arrays and total swap count.
\begin{itemize}
    \item \textbf{Pass 1:}
    \begin{itemize}
        \item Swap $7, 3 \implies [3, 7, 5, 2]$ (Swap 1)
        \item Swap $7, 5 \implies [3, 5, 7, 2]$ (Swap 2)
        \item Swap $7, 2 \implies [3, 5, 2, 7]$ (Swap 3)
    \end{itemize}
    \item \textbf{Pass 2:}
    \begin{itemize}
        \item Compare $3, 5 \implies$ no swap
        \item Swap $5, 2 \implies [3, 2, 5, 7]$ (Swap 4)
    \end{itemize}
    \item \textbf{Pass 3:}
    \begin{itemize}
        \item Swap $3, 2 \implies [2, 3, 5, 7]$ (Swap 5)
    \end{itemize}
\end{itemize}
Total Swaps: $\boxed{3 + 1 + 1 = 5 \text{ swaps}}$.

\subsection{Problem 2: Count Inversions in an Array}
An \textbf{inversion} is a pair $(i, j)$ such that $i < j$ and $A[i] > A[j]$.
\begin{itemize}
    \item For $A = [3, 1, 2]$: Inversions are $(3, 1)$ and $(3, 2) \implies \boxed{\text{Inversion Count} = 2}$.
    \item Maximum inversions in an array of size $n$:
    \[
    \boxed{I_{\max} = \frac{n(n-1)}{2}} \quad (\text{occurs on strictly reverse-sorted arrays})
    \]
    For $n = 5$ ($[5, 4, 3, 2, 1]$): $I_{\max} = \frac{5 \times 4}{2} = \boxed{10 \text{ inversions}}$.
\end{itemize}

\subsection{Problem 3: Identify the Algorithm from Its Properties}
\textbf{Question:} An algorithm runs in $O(n \log n)$ worst-case time, operates in $O(1)$ auxiliary memory, and is not stable. Which algorithm is it?\\
\textbf{Answer: Heap Sort.}
\begin{itemize}
    \item Merge Sort requires $O(n)$ space.
    \item Quick Sort has $O(n^2)$ worst case and $O(\log n)$ stack space.
    \item Insertion Sort has $O(n^2)$ worst case.
\end{itemize}

\subsection{Problem 4: Strategic Algorithm Selection}

\begin{table}[htbp]
\centering
\small
\begin{tabularx}{\textwidth}{p{6.5cm} p{3.5cm} X}
\toprule
\textbf{Real-World Operational Scenario} & \textbf{Optimal Algorithm} & \textbf{Justification} \\
\midrule
Tiny or nearly sorted array & \textbf{Insertion Sort} & Adaptive $O(n)$ time, minimal constant overhead \\
Large array needing $O(n \log n)$ guarantee and $O(1)$ space & \textbf{Heap Sort} & Worst-case optimal and in-place \\
Large dataset requiring stable sorting & \textbf{Merge Sort} & Stable and predictable $O(n \log n)$ \\
Bounded integers in small range ($k \approx n$) & \textbf{Counting Sort} & Linear $O(n + k)$ non-comparison speed \\
Fixed-width integers / Strings by digits & \textbf{Radix Sort} & Linear $O(d(n + b))$ time \\
General in-memory sorting & \textbf{Quick Sort / Introsort} & Fastest practical cache execution \\
Dataset larger than physical RAM & \textbf{External Merge Sort} & Optimized sequential disk block transfers \\
Partially ordered real-world workloads & \textbf{TimSort} & Exploits natural ordered runs \\
Strictly minimize physical memory write operations & \textbf{Selection Sort} & Performs at most $n - 1$ swaps \\
\bottomrule
\end{tabularx}
\caption{Strategic sorting algorithm selection chart.}
\end{table}

\subsection{Problem 5: Determine Sorting Stability}
\textbf{Question:} Input records: $(\text{Rahim}, 80), (\text{Karim}, 90), (\text{Salma}, 80)$.\\
Sorted ascending by marks:
\[
\text{Option A: } (\text{Rahim}, 80), (\text{Salma}, 80), (\text{Karim}, 90)
\]
\[
\text{Option B: } (\text{Salma}, 80), (\text{Rahim}, 80), (\text{Karim}, 90)
\]
\textbf{Answer: Option A is stable.} Rahim appeared before Salma in the original input; preserving this relative order confirms stability.

\subsection{Problem 6: C++ Standard Library Guarantees}
\begin{itemize}
    \item \texttt{std::sort()}: Guarantees $\mathbf{O(n \log n)}$ comparisons in the worst case (C++11 standard). Typically implemented via \textbf{Introsort}. It is \textbf{unstable}.
    \item \texttt{std::stable\_sort()}: Guarantees \textbf{stability} in $O(n \log n)$ time with auxiliary memory. Implemented via \textbf{Merge Sort / TimSort}.
\end{itemize}

\newpage
\begin{center}
\vspace{0.5em}
{\color{primary}\rule{0.45\textwidth}{1.2pt}}\\[4pt]
{\Large\bfseries\color{primary} Section X: Examination Traps and Quick Revision}\\[2pt]
{\color{primary}\rule{0.45\textwidth}{1.2pt}}
\vspace{0.5em}
\end{center}

\section{Common Sorting Examination Traps}

\begin{enumerate}
    \item \textbf{Bubble Sort Best Case:} Only $O(n)$ when the early-termination flag is implemented. Unoptimized Bubble Sort is always $\Theta(n^2)$.
    \item \textbf{Insertion Sort Comparison Operator:} Must use \texttt{A[j] > key} for stability. Using \texttt{>=} ruins stability.
    \item \textbf{Merge Sort Memory:} Conventional array Merge Sort is \textbf{not in-place}; it requires $O(n)$ buffer storage.
    \item \textbf{Quick Sort Worst Case:} Conventional Quick Sort is $\mathbf{O(n^2)}$, not $O(n \log n)$. It occurs on already sorted input with first/last pivot.
    \item \textbf{Heap Sort Stability:} Standard Heap Sort is \textbf{unstable} due to long-distance swaps between root and leaf nodes.
    \item \textbf{Selection Sort Comparisons:} Selection Sort \textbf{always} performs $\frac{n(n-1)}{2}$ comparisons ($\Theta(n^2)$), even if the array is already sorted.
    \item \textbf{Counting Sort Scope:} Runs in $O(n + k)$, not unconditional $O(n)$. It is only linear when range $k \leq O(n)$.
    \item \textbf{Radix Sort Prerequisite:} LSD Radix Sort \textbf{strictly requires} a stable inner sub-sort to preserve prior digit orders.
    \item \textbf{Shell Sort Complexity:} Cannot be stated without specifying the gap sequence! Basic halving sequence is $O(n^2)$.
    \item \textbf{TimSort vs Introsort:} TimSort is \textbf{stable} (Python); Introsort is \textbf{unstable} (C++ \texttt{std::sort}).
\end{enumerate}

\section{Final Quick-Revision Sheet}

\begin{table}[htbp]
\centering
\small
\begin{tabularx}{\textwidth}{p{3.5cm} p{2cm} p{2cm} p{2cm} c X}
\toprule
\textbf{Algorithm} & \textbf{Best Time} & \textbf{Average Time} & \textbf{Worst Time} & \textbf{Stable?} & \textbf{Auxiliary Space} \\
\midrule
Bubble Sort (Opt)  & $O(n)$ & $O(n^2)$ & $O(n^2)$ & Yes & $O(1)$ \\
Insertion Sort     & $O(n)$ & $O(n^2)$ & $O(n^2)$ & Yes & $O(1)$ \\
Selection Sort     & $O(n^2)$ & $O(n^2)$ & $O(n^2)$ & No  & $O(1)$ \\
Merge Sort         & $O(n \log n)$ & $O(n \log n)$ & $O(n \log n)$ & Yes & $O(n)$ \\
Quick Sort         & $O(n \log n)$ & $O(n \log n)$ & $O(n^2)$ & No  & $O(\log n)$ \\
Heap Sort          & $O(n \log n)$ & $O(n \log n)$ & $O(n \log n)$ & No  & $O(1)$ \\
Counting Sort      & $O(n + k)$ & $O(n + k)$ & $O(n + k)$ & Yes & $O(n + k)$ \\
Radix Sort         & $O(d(n + b))$ & $O(d(n + b))$ & $O(d(n + b))$ & Yes & $O(n + b)$ \\
TimSort            & $O(n)$ & $O(n \log n)$ & $O(n \log n)$ & Yes & $O(n)$ \\
Introsort          & $O(n \log n)$ & $O(n \log n)$ & $O(n \log n)$ & No  & $O(\log n)$ \\
\bottomrule
\end{tabularx}
\caption{Rapid-fire revision summary of key sorting algorithms.}
\end{table}

\begin{quote}
\textbf{Final Advice for the Examination Hall:} For each algorithm, ensure you can seamlessly:
\begin{enumerate}
    \item \textbf{Trace:} Execute 2--3 passes by hand on a 5-element array.
    \item \textbf{Classify:} State time complexity (best, avg, worst), space complexity, stability, and in-place status.
    \item \textbf{Justify:} Explain \textit{why} the algorithm achieves that complexity (e.g., recurrence relation, tree height, or inversion elimination).
\end{enumerate}
\end{quote}
